%++++++++++++++++++++++++++++++++++++++++++++++++++++++++++++++++++++++++++++++
% Identificacoes (se o trabalho for em inglês, insira os dados em inglês)
% Para entradas abreviadas de Professora (Profa.) em português escreva: Prof$^\textnormal{a}$.
%++++++++++++++++++++++++++++++++++++++++++++++++++++++++++++++++++++++++++++++
\course{Engenharia Aeronáutica e Mecânica} % Programa de PG ou Curso de Graduação
\area{Propulsão Aeroespacial e Energia} % Área de concentração na PG (Não utilizado no caso de TG)

% Autor do trabalho: Nome Sobrenome
\authorgender{masc} %sexo: masc ou fem
\author{\textbf{EDEN DENIS} Ferreira da Silva Lopes}{Santos}
\itaauthoraddress{Rua João Fonseca dos Santos, 20, Apto. 93}{12.230-088}{São José dos Campos--SP}

% Titulo da Tese/Dissertação
\title{Aplicação de Inteligência Artificial (AI) em Banco de Dados de Emissão de Motores}

% Orientador
\advisorgender{masc} % masc ou fem
\advisor{Prof.~Dr.}{Pedro Teixeira \textbf{LACAVA}}{ITA}

% Coorientador (Caso não haja coorientador, colocar ambas as variáveis \coadvisorgender e \coadvisor comentadas, com um % na frente)
\coadvisorgender{masc}
\coadvisor{}{}  % {Prof.~Dr.}{Carlos Alberto Rocha \textbf{PIMENTEL}}{ITA}

% Pró-reitor da Pós-graduação
\bossgender{fem}
\boss{Prof.~Dra.}{Emília \textbf{VILLANI}}

%Coordenador do curso no caso de TG
\bosscoursegender{masc}
\bosscourse{Prof.~Dr.}{ - } % masc ou fem

% Palavras-Chaves informadas pela Biblioteca -> utilizada na CIP
\kwcip{Engine taxonomy}
\kwcip{Emission classification}
\kwcip{Inteligência Artificial}
\kwcip{Artificial Intelligence}
\kwcip{Python}

% membros da banca examinadora
\examiner{Prof. Dr.}{Pedro Teixeira \textbf{LACAVA}}{President}{ITA}
\examiner{Prof. Dr.}{}{Internal Member}{ITA} % Vinícius \textbf{MALATESTA}}{Internal Member}{ITA}
\examiner{Prof. Dr.}{}{Internal Member}{ITA} % {Elisan dos Santos \textbf{MAGALHÃES}}{Internal Member}{ITA}
\examiner{Prof. Dr.}{}{Internal Member}{ITA} % {Carlos Alberto Rocha \textbf{PIMENTEL}}{External Member}{UFABC}
\examiner{Prof. Dr.}{}{Internal Member}{ITA} % {\textbf{KARL} Peter Burr}{External Member}{UFABC}

% Data da defesa (mês em maiúsculo, se trabalho em inglês, e minúsculo se trabalho em português) 
\date{14}{JULHO}{2025}

% Número CDU - (somente para TG)
\cdu{621.38}

% Folha de Registro do Documento
\newcommand{\advisorTitle}{Prof.~Dr.}
\newcommand{\advisorName}{Pedro Teixeira \textbf{LACAVA}}

\newcommand{\coadvisorTitle}{Prof.~Dr.}
\newcommand{\coadvisorName}{Carlos Alberto Rocha \textbf{PIMENTEL}}

\newcommand{\dateDay}{14}
\newcommand{\dateMonth}{JULHO}
\newcommand{\dateYear}{2025}

% Data de publicação e defesa (usada na folha de registro)
\newcommand{\frdDataPublicacao}{04 de agosto de 2025}
\newcommand{\frdNumeroRegistro}{DCTA/ITA/TD-029/2020}
\newcommand{\frdPalavrasAutor}{Big Data; Engine Emission; Emission Databank; Clustering; Artificial Intelligence}
\newcommand{\frdPalavrasResultado}{Taxonomia; Emissão de Motores; Planejamento de Dados; Inteligência Artificial; Processamento de Dados}
\newcommand{\frdApresentacao}{
	ITA, São José dos Campos. Curso de Mestrado. Programa de Pós-Graduação em Engenharia Aeronáutica e Mecânica. Área de Propulsão Aeroespacial e Energia.\\
	Orientador: \advisorTitle~\advisorName.\\
	Coorientador: \coadvisorTitle~\coadvisorName.\\
	Defesa em \dateDay/\dateMonth/\dateYear. Publicada em \frdDataPublicacao.
}

\newcommand{\PorcentagemDaLargura}{0.80}